%% SCRAP: papers/EXECUTIVE_SUMMARY
%% SOURCE: docs/working/papers/EXECUTIVE_SUMMARY.md
%% STATUS: CURRENT
%% FITS: ssrn/intro, vol3-research/intro
%% EDITORIAL: lifted — prose rewritten to press voice

\section{Executive Summary}

\subsection{The Problem}

Modern runtime systems face a persistent conflict between two design goals.
Adaptive systems---contemporary web browsers, mobile application runtimes,
JIT-compiled virtual machines---learn which code paths run most often and
reorganize themselves to accelerate those paths. They improve over time, but
their internal decisions vary from run to run; two executions of the same
program on the same machine may produce different optimization states. This
non-determinism is acceptable for web browsers but disqualifying for
safety-critical software, financial systems requiring auditable behavior, or
real-time systems with timing guarantees.

Deterministic systems---firmware for medical devices, aircraft flight
controllers, certified real-time kernels---run the same way on every
invocation. They can be formally verified, exhaustively tested, and
certified by standards bodies. But they cannot adapt: the optimization
configuration frozen at compile time may be far from optimal for the actual
workload seen at runtime.

StarForth resolves this conflict. The system is simultaneously adaptive
and completely predictable at the level of its algorithmic decisions.

\subsection{What Was Built}

StarForth is a FORTH-79 compliant virtual machine with a physics-driven
adaptive runtime. Using execution frequency as a proxy for thermal energy
(a deliberate conceptual metaphor, not a physics claim), the system tracks
how often each dictionary entry executes, applies exponential decay to reduce
the weight of stale executions, and maintains a small cache of the most
frequently executed words for accelerated lookup. Every decision in this
pipeline is deterministic: given the same execution sequence, the system
produces the same cache configuration on every run.

Seven feedback loops coordinate the adaptive behavior:

\begin{itemize}
  \item Execution heat tracking and hot-words caching (Loop~1)
  \item Rolling window of execution history (Loop~2)
  \item Linear decay of quiescent entries (Loop~3)
  \item Word-to-word transition prediction (Loop~4)
  \item Variance-based window width inference via Levene's test (Loop~5)
  \item Decay slope inference via exponential regression (Loop~6)
  \item Adaptive heartbeat coordination (Loop~7)
\end{itemize}

All loops are independently togglable and formally characterized. Fixed-point
arithmetic (\Qtype\ throughout) eliminates IEEE-754 non-determinism across
architectures.

\subsection{Experimental Results}

A three-configuration design of experiments ran 90 trials (30 runs per
configuration):

\begin{itemize}
  \item \textbf{C\_NONE} (baseline): all adaptive loops disabled
  \item \textbf{C\_CACHE}: hot-words cache enabled, inference disabled
  \item \textbf{C\_FULL}: all seven loops active
\end{itemize}

\paragraph{Result 1: algorithmic determinism.}
Across all 30 runs of C\_FULL, the cache hit rate was 17.39\% with a
coefficient of variation (CV) of exactly 0.00\%. The F-test for variance
homogeneity yields $F(29,29) \to \infty$, $p < 10^{-30}$. Even under
intentional thermal stress (sustained CPU load), the algorithmic decisions
did not change; only wall-clock runtime increased.

\paragraph{Result 2: adaptive convergence.}
C\_FULL improved from a mean of 10.20\,ms in early runs (1--15) to 7.61\,ms
in late runs (16--30), a statistically significant 25.4\% improvement
($t = 4.23$, $p < 0.001$, Cohen's $d \approx 5.08$). The baseline
configuration (C\_NONE) showed slight degradation over the same window;
C\_CACHE showed no meaningful improvement. Only the fully adaptive
configuration converged.

\paragraph{Result 3: variance decomposition.}
Runtime exhibits 60--70\% CV due to OS scheduler noise, thermal variation,
and cache-line effects. Cache decisions exhibit 0.00\% CV. The two components
are statistically independent (Pearson $r = 0.03$, $p = 0.87$). The
adaptive algorithm is perfectly stable; the environment adds noise on top
of a deterministic foundation.

\subsection{Implications}

\paragraph{Verifiable adaptive systems.}
A system whose adaptive decisions are deterministic can, in principle, be
formally verified. The optimization state reachable from a given workload is
predictable, auditable, and reproducible---properties required for
safety-critical certification.

\paragraph{Reproducible performance characterization.}
Benchmark results become portable: different researchers running the same
workload reach the same adaptive steady state, enabling apples-to-apples
comparison of optimization strategies.

\paragraph{Adaptation without non-determinism.}
The results demonstrate that statistical inference can guide optimization
without introducing randomness. The system adapts through arithmetic and
counting, not machine learning or stochastic search.

\subsection{Scope and Limitations}

This work does not claim that StarForth is the fastest virtual machine, that
its approach is superior to JIT compilation for all use cases, or that the
thermodynamic framework is a physical theory. The 0.00\% algorithmic CV
applies to cache decisions, not to total system runtime. The system is
validated on CPU-bound, deterministic FORTH programs; I/O-bound workloads,
programs with random control flow, and very short processes fall outside the
validated scope. Fifteen failure modes are documented in the companion
negative-results section.

\subsection{Reproducibility}

Full source code, raw experimental data from all 90 runs, and a
step-by-step reproduction protocol are publicly available. A Docker container
provides bit-for-bit exact reproduction against a pinned environment. The
authoritative reproduction command is:

\begin{lstlisting}[language=bash]
git clone https://github.com/rajames440/StarForth.git
cd StarForth
make fastest
./build/amd64/fastest/starforth --doe
\end{lstlisting}

Deviations from the expected 0.00\% algorithmic CV should be reported as bugs.

\subsection{Summary}

\begin{itemize}
  \item An adaptive virtual machine can achieve 0.00\% algorithmic variance
    across 90 experimental runs.
  \item Adaptation and determinism are not mutually exclusive; statistical
    inference drives optimization without introducing non-determinism.
  \item Performance improves measurably (25.4\%) as the system converges to
    steady state.
  \item All data, code, and reproduction instructions are publicly available
    for independent verification.
\end{itemize}
